% Compilen desde esta carpeta o desde la raíz del repositorio.
\makeatletter
\def\input@path{{../../preamble/}{book2026-2027/preamble/}}
\makeatother
\input{preamble.tex}

\usepackage[spanish]{babel}
\usepackage{lmodern}
\usepackage{titlesec}
\geometry{a4paper,margin=1.6cm}
\linespread{1.10}

\definecolor{HeadingColor}{HTML}{243B53}
\definecolor{InstructionColor}{HTML}{704214}
\definecolor{ConceptColor}{HTML}{205C50}
\definecolor{LinkColor}{HTML}{0047AB}
\hypersetup{allcolors=LinkColor}
\newcommand{\guidelink}[2]{\href{#1}{\textbf{\underline{#2}}}}
\newcommand{\instruction}[1]{\textcolor{InstructionColor}{\textbf{#1}}}
\newcommand{\concept}[1]{\textcolor{ConceptColor}{\textbf{#1}}}
\newcommand{\checkbox}{%
    \begingroup
    \setlength{\fboxrule}{0.4pt}%
    \setlength{\fboxsep}{2pt}%
    \fbox{\rule{0pt}{0.8ex}\rule{0.8ex}{0pt}}%
    \endgroup}

\setlength{\parskip}{4pt}
\setlength{\tabcolsep}{5pt}
\renewcommand{\arraystretch}{1.15}
\setlist{leftmargin=0.65cm,topsep=4pt,partopsep=0pt,parsep=0pt,itemsep=3pt}
\titleformat{\section}{\large\bfseries\color{HeadingColor}}{}{0pt}{}
\titlespacing*{\section}{0pt}{10pt}{3pt}
\newcolumntype{Y}{>{\raggedright\arraybackslash}X}

\hypersetup{pdftitle={C3--C4. Representen y construyan la primera versión},pdfsubject={Emprendimiento e Innovación, Grado 10, Periodo 2}}

\begin{document}

\begin{center}
{\Large\bfseries\color{HeadingColor} C3--C4. Representen y construyan la primera versión\par}
\vspace{3pt}
Emprendimiento e Innovación \enspace|\enspace Grado 10 \enspace|\enspace Periodo 2\\
Guía 2 de 5
\end{center}

\section*{Objetivo y punto de partida}
\concept{C3.} Describe la propuesta mediante bocetos, diagramas, guiones, mapas de servicio o esquemas funcionales.\\
\concept{C4.} Construye un prototipo de baja fidelidad físico, digital, visual, comunicativo o procedimental.

Recuperen D1 y los H de P1, y la alternativa elegida y sus supuestos (P2, C1--C2). Elijan qué función o interacción conviene explorar primero y qué quieren aprender de ella.

\section*{Documento acumulativo e integridad de la evidencia}
Esta guía es de consulta. Continúen el \textbf{mismo documento del proyecto en LaTeX} en \guidelink{https://www.overleaf.com/}{Overleaf} o en su editor habitual. Actualicen la fecha y entreguen el PDF completo. Usen \textbf{P1} para el Periodo 1 y \textbf{P2} para el Periodo 2 al citar criterios. Enlacen los registros por código o referencia interna y conserven sus fuentes en la sección de origen.

Distingan \concept{evidencia} (lo observado o expresado), \concept{conclusión} (lo que infieren), \concept{supuesto} (lo pendiente de verificar) e \concept{idea} (una propuesta). Registren fuentes y fechas reales. Identifiquen toda simulación autorizada por el docente. \instruction{Si usan IA, limítenla al código LaTeX, al formato y a organizar contenido propio.} Evidencias, conclusiones, ideas, decisiones y aportes deben provenir de su trabajo.

\section*{C3 de P2. Hagan visible cómo funcionaría la propuesta}
\begin{enumerate}[series=pasos]
\item \textbf{Elijan qué representar.} Retomen un supuesto o una condición de la decisión. Delimiten una acción concreta que realizará el usuario e indiquen qué debería ocurrir después.
\item \textbf{Usen un formato adecuado al proyecto.} Elijan un boceto anotado para un objeto, pantallas o un diagrama para una interacción digital, un guion para una pieza comunicativa, o un mapa de servicio para un proceso. Identifiquen al usuario, el punto de partida, los pasos y la respuesta esperada. Señalen qué parte convertirán en prototipo.
\item \textbf{Agreguen la representación al documento.} Asígnenle el código R1, pónganle un título y vincúlenla con D1, la alternativa ALT elegida y el hallazgo H o supuesto que explica el recorte.
\end{enumerate}

\textbf{Ejemplo ficticio.} En un mapa de servicio para recoger pedidos, marquen dónde una persona busca su número, dónde se ubica y qué indicación recibe. Para ensayar esa interacción, pueden simular la señal y la entrega con tarjetas y dos puestos.

\concept{Evidencia de C3.} Una representación anotada que permite seguir la acción del usuario e identificar la parte que se llevará al prototipo.

\newpage
\section*{C4 de P2. Planeen y construyan el prototipo V1}
\begin{enumerate}[resume=pasos]
\item \textbf{Acuerden un plan breve antes de construir.} Abran el registro PL1. Dividan la construcción en tareas cortas, asignen responsables y roles, confirmen recursos y acuerden plazos. Anoten qué permiso, material o decisión queda pendiente. Conserven esta primera planificación para revisarla en P2, C6.
\end{enumerate}

{\small
\begin{tabularx}{\linewidth}{|Y|Y|Y|Y|}
\hline
\textbf{Tarea} & \textbf{Responsable y rol} & \textbf{Recursos y acceso} & \textbf{Plazo y pendiente} \\
\hline
Parte de R1 que prepararán. & Quién la prepara y qué hace. & Material, herramienta o espacio disponible. & Cuándo debe quedar lista y qué falta confirmar. \\
\hline
\end{tabularx}\par
}

\begin{enumerate}[resume=pasos]
\item \textbf{Construyan y recorran V1.} Usen materiales sencillos: papel, cartón, pantallas enlazadas, tarjetas, un guion representado u otro medio apropiado. Hagan un recorrido interno en el que una persona ejecute la acción y otra simule la respuesta si hace falta. Identifiquen qué partes funcionan, cuáles simulan y qué queda por construir.
\item \textbf{Incorporen V1 al documento.} Incluyan fotos, capturas o una secuencia con pies que muestren la acción y su respuesta. Identifiquen la versión y la fecha, enlacen R1 y PL1, y anoten qué aprendieron del recorrido interno. Actualicen la bitácora con aportes, compromisos y acuerdos de la sesión.
\end{enumerate}

\concept{Evidencia de C4.} Una primera versión de baja fidelidad que permite ensayar la acción elegida, con un registro visual o secuencial que muestra cómo se usa y qué se simula.

\section*{Resultado y continuidad}
Agreguen R1, PL1 y V1. En P2, C5--C6 usarán V1 para aplicar decisiones de diseño y retomarán el mismo PL1 para preparar la siguiente versión y su prueba.

\section*{Verificación final}
\begin{itemize}[label=\checkbox]
\item R1 muestra quién actúa, qué hace y qué respuesta espera.
\item PL1 quedó registrado antes de construir, con tareas, responsables, recursos y plazos.
\item V1 permite recorrer la interacción elegida y documentamos sus límites.
\item Actualizamos la bitácora y compilamos el PDF acumulativo con sus referencias.
\end{itemize}

\end{document}
